\documentclass[../../main.tex]{subfiles}

\begin{document}

\section{对偶空间和弱序列收敛}\label{section:ch16-dual-weak}

对 $E$ 是实数的可测集, $1\leqslant p<\infty$, 而 $q$ 是 $p$ 的共轭, 关于 $L^{p}(E)$ 的 Riesz 表示定理直接描述了 $L^{q}(E)$ 映上 $[L^{p}(E)]^{*}$ 的线性等距. 该定理的 $p=2$ 情形可推广到一般的 Hilbert 空间.

\begin{theorem}[Riesz-Fr\'echet 表示定理]\label{theorem:ch16-riesz-frechet}
令 $H$ 为 Hilbert 空间. 定义算子 $T:H\to H^{*}$ 为对每个 $h\in H$ 赋予线性泛函 $T(h):H\to\mathbf{R}$, 它定义为
\begin{align}
\text{对所有 }u\in H,\quad T(h)[u]=\langle h,u\rangle.\label{align:ch16-riesz-1}
\end{align}
则 $T$ 是 $H$ 映上 $H^{*}$ 的线性等距.
\end{theorem}

\begin{proof}
令 $h$ 属于 $H$. 则 $T(h)$ 是线性的. 我们从 Cauchy-Schwarz 不等式推出泛函 $T(h):H\to\mathbf{R}$ 是有界的且 $\|T(h)\|\leqslant\|h\|$. 但若 $h\neq 0$, 则 $T(h)[h/\|h\|]=\|h\|$. 因此 $\|T(h)\|=\|h\|$. 因此 $T$ 是等距. 显然 $T$ 是线性的. 剩下来要证明 $T(H)=H^{*}$. 令 $\psi_{0}\neq 0$ 属于 $H^{*}$. 由于 $\psi_{0}$ 是连续的, 它的核是 $H$ 的闭真子空间. 根据定理\cref{theorem:ch16-orthogonal-decomposition}, 由于 $\ker\psi_{0}\neq H$, 我们可以选取单位向量 $h^{*}\in H$, 它正交于 $\ker\psi_{0}$. 定义 $h_{0}=\psi_{0}(h.)h_{.}$. 我们宣称 $T(h_{0})=\psi_{0}$. 的确, 对 $h\in H$,
\begin{align*}
h-\frac{\psi_{0}(h)}{\psi_{0}(h.)}h_{.}\in\ker\psi_{0},\quad\text{因此 }\langle h-\frac{\psi_{0}(h)}{\psi_{0}(h.)}h_{.},\ h.\rangle=0.
\end{align*}
所以 $\psi_{0}(h)=\langle h_{0},\ h\rangle=T(h_{0})[h]$. 因此 $T(H)=H^{*}$.
\end{proof}

如同在一般的 Banach 空间的情形, 对 Hilbert 空间 $H$ 中的序列, 我们写 $\{u_{n}\}\to u$ 意味着 $H$ 中的序列弱收敛于 $u\in H$. 根据 Riesz-Fr\'echet 表示定理, 在 $H$ 中,
\begin{align*}
\{u_{n}\}\to u\text{ 当且仅当对所有 }h\in H,\quad\lim_{n\to\infty}\langle h,u_{n}\rangle=\langle h,u\rangle.
\end{align*}

\begin{theorem}\label{theorem:ch16-weak-seq}
Hilbert 空间 $H$ 中的每个有界序列有弱收敛的子序列.
\end{theorem}

\begin{proof}
令 $\{h_{n}\}$ 为 $H$ 中的有界序列. 定义 $H_{0}$ 为 $\{h_{n}\}$ 的闭线性扩张. 则 $H_{0}$ 是可分的. 对每个自然数 $n$, 定义 $\psi_{n}\in[H_{0}]^{*}$ 为
\begin{align*}
\text{对所有 }h\in H_{0},\quad\psi_{n}(h)=\langle h_{n},h\rangle.
\end{align*}
由于 $\{h_{n}\}$ 是有界的, 我们从 Cauchy-Schwarz 不等式推出 $\{\psi_{n}\}$ 也是有界的. 则 $\{\psi_{n}\}$ 是可分赋范线性空间 $H_{0}$ 上的有界线性泛函的有界序列. Helly 定理告诉我们存在 $\{\psi_{n}\}$ 的子序列 $\{\psi_{n_{k}}\}$ 逐点收敛于 $\psi_{0}\in[H_{0}]^{*}$. 根据 Riesz-Fr\'echet 表示定理, 存在向量 $h_{0}\in H_{0}$, 使得 $\psi_{0}=T(h_{0})$. 因此,
\begin{align*}
\text{对所有 }h\in H_{0},\quad\lim_{k\to\infty}\langle h_{n_{k}},h\rangle=\langle h_{0},h\rangle.
\end{align*}
令 $P$ 为 $H$ 映上 $H_{0}$ 的正交投影. 对每个指标 $k$, 由于 $(\operatorname{Id}-P)[H]=P(H)^{\perp}$,
\begin{align*}
\text{对所有 }h\in H,\quad\langle h_{n_{k}},(\operatorname{Id}-P)[h]\rangle=\langle h_{0},(\operatorname{Id}-P)[h]\rangle=0.
\end{align*}
因此,
\begin{align*}
\text{对所有 }h\in H,\quad\lim_{k\to\infty}\langle h_{n_{k}},h\rangle=\langle h_{0},h\rangle.
\end{align*}
因此 $\{h_{n_{k}}\}$ 在 $H$ 中弱收敛于 $h_{0}$.
\end{proof}

在以下命题中我们收集了关于弱收敛序列的一些性质, 我们先前已对一般的 Banach 空间证明了这些性质, 但因为有了 Riesz-Fr\'echet 表示定理, 在 Hilbert 空间的情形下证明它们要简单得多(见习题\cref{exercise:ch16-17}).

\begin{proposition}\label{proposition:ch16-weak-properties}
在 Hilbert 空间 $H$ 中, 弱意义下 $\{u_{n}\}\to u$. 则 $\{u_{n}\}$ 是有界的且
\begin{align*}
\|u\|\leqslant\liminf\|u_{n}\|.
\end{align*}
此外, 若在 $H$ 中, 强意义下 $\{v_{n}\}\to v$, 则
\begin{align}
\lim_{n\to\infty}\langle u_{n},v_{n}\rangle=\langle u,v\rangle.\label{align:ch16-weak-inner}
\end{align}
\end{proposition}

以下两个命题描述了 Hilbert 空间中的弱收敛序列的性质, 在第 8 章我们已注意到它们对 $L^{p}(E)$ 空间成立, 其中 $E$ 是实数的可测集而 $1<p<\infty$, 但对一般的 Banach 空间这些性质不成立.

\begin{theorem}[Banach-Saks 定理]\label{theorem:ch16-banach-saks}
在 Hilbert 空间 $H$ 中弱意义下 $\{u_{n}\}\to u$. 则存在 $\{u_{n}\}$ 的子序列 $\{u_{n_{k}}\}$ 使得
\begin{align}
\text{在 }H\text{ 中强意义下 }\lim_{k\to\infty}\frac{u_{n_{1}}+\cdots+u_{n_{k}}}{k}=u.\label{align:ch16-banach-saks-1}
\end{align}
\end{theorem}

\begin{proof}
用 $u_{n}-u$ 代替每个 $u_{n}$, 我们可以假定 $u=0$. 由于弱收敛的序列是有界的, 我们可以选取 $M>0$ 使得
\begin{align*}
\text{对所有 }n,\quad\|u_{n}\|^{2}\leqslant M.
\end{align*}
我们将归纳地选取 $\{u_{n}\}$ 的子序列 $\{u_{n_{k}}\}$, 具有性质: 对所有 $k$,
\begin{align}
\|u_{n_{1}}+\cdots+u_{n_{k}}\|^{2}\leqslant(2+M)k.\label{align:ch16-banach-saks-2}
\end{align}
对这样的序列,
\begin{align}
\text{对所有 }k,\quad\left\|\frac{u_{n_{1}}+\cdots+u_{n_{k}}}{k}\right\|^{2}\leqslant\frac{(2+M)}{k}.\label{align:ch16-banach-saks-3}
\end{align}
这样就结束了证明.

定义 $n_{1}=1$. 由于 $\{u_{k}\}\to 0$ 且 $u_{n_{1}}$ 属于 $H$, 我们能够选取指标 $n_{2}$ 使得 $|\langle u_{n_{1}},\ u_{n_{2}}\rangle|<1$. 假定我们已选取自然数 $n_{1}<n_{2}<\cdots<n_{k}$ 使得
\begin{align*}
\text{对 }j=1,\cdots,k,\quad\|u_{n_{1}}+\cdots+u_{n_{j}}\|^{2}\leqslant(2+M)j.
\end{align*}
由于 $\{u_{n}\}\to 0$ 且 $u_{n_{1}}+\cdots+u_{n_{k}}$ 属于 $H$, 我们可以选取 $n_{k+1}>n_{k}$ 使得
\begin{align*}
|\langle u_{n_{1}}+\cdots+u_{n_{k}},u_{n_{k+1}}\rangle|\leqslant 1.
\end{align*}
然而,
\begin{align*}
\|u_{n_{1}}+\cdots+u_{n_{k}}+u_{n_{k+1}}\|^{2}=\|u_{n_{1}}+\cdots+u_{n_{k}}\|^{2}+2\langle u_{n_{1}}+\cdots+u_{n_{k}},u_{n_{k+1}}\rangle+\|u_{n_{k+1}}\|^{2}.
\end{align*}
因此,
\begin{align*}
\|u_{n_{1}}+\cdots+u_{n_{k+1}}\|^{2}\leqslant(2+M)k+2+M=(2+M)(k+1).
\end{align*}
我们已选取子序列使得\eqref{align:ch16-banach-saks-2}成立.
\end{proof}

\begin{theorem}[Radon-Riesz 定理]\label{theorem:ch16-radon-riesz}
在 Hilbert 空间 $H$ 中弱意义下 $\{u_{k}\}\to u$. 则在 $H$ 中强意义下 $\{u_{k}\}\to u$ 当且仅当 $\lim_{n\to\infty}\|u_{n}\|=\|u\|$.
\end{theorem}

\begin{proof}
由于范数在 $H$ 上关于强拓扑连续, 若强意义下 $\{u_{k}\}\to u$, 则 $\lim_{n\to\infty}\|u_{n}\|=\|u\|$. 另一方面, 若 $\lim_{n\to\infty}\|u_{n}\|=\|u\|$, 则由于
\begin{align*}
\text{对所有 }n,\quad\|u_{n}-u\|^{2}=\|u_{n}\|^{2}-2\langle u_{n},u\rangle+\|u\|^{2},
\end{align*}
弱收敛序列 $\{u_{k}\}$ 是强收敛的.
\end{proof}

\begin{theorem}\label{theorem:ch16-reflexive}
令 $H$ 为 Hilbert 空间. 则 $H$ 是自反的. 因此 $H$ 的每个非空强闭有界凸子集 $K$ 是弱紧的, 且因此是它的极值点的强闭凸包.
\end{theorem}

\begin{proof}
为证明自反性, 必须证明自然嵌入 $J:H\to[H^{*}]^{*}$ 是映上的. 令 $\psi:H^{*}\to\mathbf{R}$ 为有界线性泛函. 令 $T:H\to[H]^{*}$ 为 Riesz-Fr\'echet 表示定理描述的同构. 则 $\psi\circ T:H\to\mathbf{R}$ 是有界线性算子的复合, 故是有界的. Riesz-Fr\'echet 表示定理告诉我们存在 $h_{0}\in H$ 使得 $\psi\circ T=T(h_{0})$. 因此,
\begin{align*}
\text{对所有 }h\in H,\quad\psi(T(h))=T(h_{0})[h]=T(h)[h_{0}]=J(h_{0})[T(h)].
\end{align*}
由于 $T(H)=H^{*}$, $\psi=J(h_{0})$. 因此 $H$ 是自反的. 我们从 Kakutani 定理与 Mazur 定理推出 $H$ 的每个非空强闭有界凸子集 $K$ 是弱紧的. 因此, 根据 Krein-Milman 定理与 Mazur 定理的另外一个应用, 这样的集合 $K$ 是它的极值点的强闭凸包.
\end{proof}

\begin{exercise}\label{exercise:ch16-16}
证明 $\ell^{1},\ \ell^{\infty},\ L^{1}[a,\ b],\ L^{\infty}[a,\ b]$ 都不是 Hilbert 空间.
\end{exercise}

\begin{exercise}\label{exercise:ch16-17}
证明命题\cref{proposition:ch16-weak-properties}.
\end{exercise}

\begin{exercise}\label{exercise:ch16-18}
令 $H$ 为内积空间. 证明由于 $H$ 是 Banach 空间 $X$ 的稠密子集, $X$ 的范数限制在由 $H$ 上的内积诱导的范数, $H$ 上的内积可延拓到 $X$ 上且诱导了 $X$ 上的范数. 因此内积空间有 Hilbert 空间完备化.
\end{exercise}

\end{document}